\section{Configura\c{c}\~ao Experimental}
\label{sec:setup}

A avaliação mede a precisão da evidência das UDVs por anotação humana (Seção~\ref{sec:setup-udv}) e
o efeito de cada elemento da simulação por perguntas de múltipla escolha
(Seção~\ref{sec:setup-simulacao}). O gabarito dessas perguntas são opiniões publicadas sobre
audiências posteriores às do perfil, ligadas aos atores pelas UDVs, e as medidas da geração aberta são
descritivas.

\subsection{Evidência das UDVs}
\label{sec:setup-udv}

% Pendente: dizer quais alternativas do vínculo por similaridade semântica foram comparadas nas audiências de validação e onde essa comparação é relatada.
A precisão da evidência é medida numa amostra cega sorteada somente nas audiências de teste, porque
as de validação foram usadas na comparação de alternativas do vínculo por similaridade semântica. A
amostra tem 127 linhas (134 UDVs, 30 audiências): 35 das 41 UDVs \texttt{quote\_found}, 65 das 277
\texttt{semantic\_match\_high} sem coincidência de prefixo curto, 15 das 22 UDVs semânticas com essa
coincidência, as 6 \texttt{semantic\_match\_weak} restantes e as 13 UDVs sem evidência. Há
coincidência de prefixo curto quando o prefixo mais longo da citação encontrado na fala tem menos de
seis palavras e ocorre pela primeira vez na sentença escolhida pelo codificador, caso que a UDV
registra como corroboração fraca (Apêndice~\ref{ap:regras-udv}). Os totais desses cinco grupos somam as 359 opiniões de teste
(Seção~\ref{sec:dados-particao}). As 13 UDVs sem evidência ocupam 6 linhas, uma por pessoa, nas quais
o anotador indica se a pessoa falou na sessão.

Nas demais linhas, o anotador vê a opinião, a sentença e o texto ao redor, sem nível, cosseno ou
marca de corroboração fraca, e julga a sentença como correta, parcial ou incorreta. Nas linhas de
citação literal e de prefixo curto, as palavras citadas reaparecem na sentença, e por isso a
distinção entre os dois níveis semânticos é a única que permanece cega. Ao menos 24 horas depois da
primeira anotação, 20 linhas voltam em outra ordem, para medir a concordância do anotador consigo
mesmo.

Os critérios, registrados antes de qualquer julgamento, exigem que o limite inferior do intervalo de
Wilson~\citep{wilson1927probable} de 95\% atinja 0,90 na precisão estrita (só ``correta'') da citação
literal. No nível \texttt{semantic\_match\_high} sem prefixo curto, o mesmo limite deve atingir 0,75
na precisão tolerante (``correta'' ou ``parcial''). Com as linhas sorteadas, isso exige 35 corretas
em 35 e ao menos 56 aceitas em 65.

\subsection{Simulação}
\label{sec:setup-simulacao}

Na simulação, a avaliação mede se o perfil, os trechos e a CFG aumentam a capacidade do modelo
simulador de reconhecer a opinião que a matéria atribuiu à pessoa numa audiência que não entrou no
perfil. Essa opinião é apresentada entre outras opiniões publicadas sobre a mesma audiência.

A fala gerada não serve como medida principal, porque o assunto da audiência, fornecido ao modelo,
pode já conter a posição central da opinião. Por exemplo, na audiência 1, do teste, o assunto
registrado é ``Acusações de censura contra Alexandre de Moraes por exigir bloqueio de contas na
rede social X''. Uma das opiniões dessa audiência é ``Acusou Alexandre de Moraes de censura ao
solicitar o bloqueio de contas na rede social X''. Na múltipla escolha, o assunto é o mesmo para
todas as opções, e o acerto ao acaso é conhecido (25\%).

\textbf{Perguntas.} Cada UDV de uma audiência de validação ou de teste cujo turno de evidência
pertence a um ator com perfil gera uma pergunta, acompanhada da data e do assunto da audiência. O
texto da pergunta é ``Uma destas afirmações foi feita por \textit{nome} nesta audiência, que não
aparece no material desta conversa. Qual é a mais provável?''. As quatro opções são a opinião da
UDV, na redação do resumo estruturado, e opiniões atribuídas pela matéria a três outros participantes
da mesma audiência, sorteados com semente fixa. As opiniões da própria pessoa e os textos iguais a
elas são excluídos, um texto atribuído a mais de um participante é incluído uma única vez, e a
pergunta com menos de três outros participantes é descartada. Como os distratores vêm da mesma
audiência, eles tratam do mesmo tema, o que reduz, sem eliminar, a possibilidade de escolher a opção
pela proximidade com o assunto.

Na pergunta, a UDV apenas liga a opinião a um registro de ator, pelo turno da evidência, e a
sentença de evidência não aparece. A validade da pergunta depende, portanto, da resolução da pessoa
na audiência e da identificação do ator entre audiências (Seção~\ref{sec:metodo-atores}); ela
independe da precisão do vínculo por similaridade semântica. Das 106 opiniões de teste ligadas aos
atores com falas de treino (Seção~\ref{sec:dados-particao}), 17 têm nível \texttt{quote\_found}, 87
\texttt{semantic\_match\_high} e 2 \texttt{semantic\_match\_weak}.

\textbf{Leitura da resposta.} A escolha é lida na distribuição de probabilidade do primeiro token
gerado, restrita aos tokens das quatro letras, e antes da execução confere-se que cada letra é um
único token. Como LLMs favorecem certas posições entre as
opções~\citep{zheng2024large,pezeshkpour2024large}, cada pergunta é apresentada nas quatro rotações
cíclicas de uma ordem sorteada. Em cada rotação, as probabilidades das letras (na condição 3, já
combinadas pela Equação~\ref{eq:cfg}) são renormalizadas entre si. A probabilidade de cada opção é a
média das rotações, a resposta é a opção de maior probabilidade, e as métricas são o acerto e a
probabilidade média da opção correta.

Como a leitura sempre escolhe uma letra e a probabilidade do primeiro token pode divergir da
resposta escrita~\citep{wang2024answer}, as condições 0 a 2 registram como diagnóstico a soma das
probabilidades das quatro letras no vocabulário inteiro. Essa soma mostra quanto da distribuição do
primeiro token cabe às letras.

\textbf{Seleção de $k$ e $\gamma$.} Os dois hiperparâmetros são escolhidos nas audiências de
validação, e antes da execução verifica-se que nenhum perfil e nenhuma fala usada na recuperação
contém audiência de validação ou de teste. A verificação não cobre o cargo, que na múltipla escolha
vem da matéria da audiência avaliada. Primeiro, $k \in \{3, 5, 10\}$ é o valor de maior acerto da
condição 2. Depois, com esse $k$, $\gamma \in \{1;\ 1{,}5;\ 2;\ 3\}$ é o de maior acerto da condição 3,
para que a diferença entre as condições 2 e 3 dependa só da CFG. Empates são resolvidos pela maior
probabilidade média da opção correta e, depois, pelo menor valor da grade. A escolha de
$\gamma = 1$, que reproduz a condição 2, indicaria que a CFG não aumentou o acerto na validação, e a
condição 3 coincidiria com a 2 no teste. No teste, as condições são medidas com o $k$ e o $\gamma$
escolhidos.

\textbf{Comparação entre condições.} A diferença de acerto entre condições vizinhas estima o efeito
do elemento acrescentado, e o intervalo de 95\% de cada diferença é obtido por \textit{bootstrap}
pareado de percentis, com 1.000 reamostragens de audiências inteiras~\citep{efron1993introduction}.
As audiências são reamostradas inteiras porque as perguntas de uma audiência compartilham assunto,
falantes e distratores. Os três intervalos não são ajustados para comparações múltiplas, e as demais
medidas são descritivas.

\textbf{Base da estimativa.} No teste, as condições 1 e 2 recebem também o pedido da base da
estimativa, sem as opções, para que a classificação não dependa dos distratores. A condição 3 herda
a classificação da 2, e as respostas sem rótulo são contadas à parte. Pelo critério registrado no
desenho da avaliação antes da execução, a classificação só é considerada informativa se o acerto na
múltipla escolha for menor nas perguntas ESPECULATIVA que nas DIRETA.

\textbf{Geração aberta.} A geração aberta, que corresponde ao uso previsto do simulador, é feita
para cada par formado por um ator com perfil e uma audiência de teste com opinião ligada a ele, e
suas medidas são descritivas. O modelo recebe a data e o assunto da audiência e gera a base da
estimativa, a fala e a justificação nas condições 1 e 2, com o $k$ escolhido na múltipla escolha.
Como o prompt da geração aberta traz os exemplos de fala, os resultados não são comparáveis aos da
múltipla escolha na mesma condição.

As falas não são comparadas às opiniões publicadas, e as medidas são as recusas, a fração de frases
sem fundamento e as respostas interrompidas pelo limite de tokens (400 para a fala e 1.500 para a
justificação). Uma frase apoiada em item do perfil é aceita sem cópia, e uma apoiada em exemplo ou
trecho exige cópia conferida. Por isso, a fração de frases sem fundamento depende do material de
cada condição e mede a conformidade da justificação com esse material, sem avaliar a correção da
fala. Também é gerada a fala com o material da condição 0, o prompt negativo da CFG, e a leitura
dessas falas conta as recusas. Essa contagem verifica, na execução completa, o resultado do teste
preliminar que motivou a retirada da CFG da geração aberta (Seção~\ref{sec:metodo-simulacao}).

\textbf{Recursos.} O modelo simulador é executado em <hardware> com <precisão numérica>, e o
codificador é o Serafim 335m, na mesma versão de pesos da UDV. A classificação da base da estimativa
tem limite de 8 tokens de saída, e o código, as configurações, os prompts e as saídas estão em
<repositório>.
